\textbf{Phần 1. Trắc nghiệm lựa chọn}
\setcounterpageref{ex}{0}

\begin{ex} %[2H7N35]
Đơn vị dùng để đo khối lượng là
\choice
{Xăng-ti-mét}
{Lít}
{\True Ki-lô-gam}
{Giờ}
\loigiai{
Đo khối lượng dùng ki-lô-gam.\\
\textbf{Đáp án của câu hỏi là C}
}
\end{ex}

\begin{ex} %[2H7N35]
Đơn vị dùng để đo dung tích (sức chứa chất lỏng) là
\choice
{Ki-lô-gam}
{\True Lít}
{Mét}
{Ngày}
\loigiai{
Đo dung tích dùng lít.\\
\textbf{Đáp án của câu hỏi là B}
}
\end{ex}

\begin{ex} %[2H7H35]
Tính: $38\text{ kg} + 27\text{ kg} = \dots$
\choice
{$64\text{ kg}$}
{$66\text{ kg}$}
{$67\text{ kg}$}
{\True $65\text{ kg}$}
\loigiai{
$38 + 27 = 65\text{ kg}$.\\
\textbf{Đáp án của câu hỏi là D}
}
\end{ex}

\begin{ex} %[2H7H35]
Tính: $72\text{ l} - 38\text{ l} = \dots$
\choice
{\True $34\text{ l}$}
{$33\text{ l}$}
{$35\text{ l}$}
{$36\text{ l}$}
\loigiai{
$72 - 38 = 34\text{ l}$.\\
\textbf{Đáp án của câu hỏi là A}
}
\end{ex}

\begin{ex} %[2H7H35]
Lúc 19 giờ 30 phút tối còn gọi là
\choice
{$6$ giờ $30$ phút tối}
{\True $7$ giờ $30$ phút tối (hoặc 7 giờ rưỡi tối)}
{$8$ giờ $30$ phút tối}
{$9$ giờ $30$ phút tối}
\loigiai{
$19 - 12 = 7$ giờ 30 phút tối.\\
\textbf{Đáp án của câu hỏi là B}
}
\end{ex}

\begin{ex} %[2H7H35]
Tháng 7 có bao nhiêu ngày?
\choice
{$30$ ngày}
{$28$ ngày}
{$29$ ngày}
{\True $31$ ngày}
\loigiai{
Tháng 7 có 31 ngày.\\
\textbf{Đáp án của câu hỏi là D}
}
\end{ex}

\begin{ex} %[2H7V35]
Bao gạo to nặng $45\text{ kg}$, bao gạo nhỏ nhẹ hơn bao gạo to $18\text{ kg}$. Bao gạo nhỏ nặng số ki-lô-gam là
\choice
{\True $27\text{ kg}$}
{$26\text{ kg}$}
{$28\text{ kg}$}
{$29\text{ kg}$}
\loigiai{
$45 - 18 = 27\text{ kg}$.\\
\textbf{Đáp án của câu hỏi là A}
}
\end{ex}

\begin{ex} %[2H7V35]
Một can chứa $25\text{ l}$ mật ong, đã bán đi $9\text{ l}$. Can còn lại số lít mật ong là
\choice
{$15\text{ l}$}
{$17\text{ l}$}
{\True $16\text{ l}$}
{$18\text{ l}$}
\loigiai{
$25 - 9 = 16\text{ l}$.\\
\textbf{Đáp án của câu hỏi là C}
}
\end{ex}

\begin{ex} %[2H7V35]
Nếu hôm nay là thứ Sáu ngày 9 thì thứ Sáu tuần sau là ngày mấy?
\choice
{$15$}
{\True $16$}
{$17$}
{$18$}
\loigiai{
$9 + 7 = 16$.\\
\textbf{Đáp án của câu hỏi là B}
}
\end{ex}

\begin{ex} %[2H7V35]
Tính giá trị biểu thức: $45\text{ kg} + 35\text{ kg} - 28\text{ kg}$.
\choice
{\True $52\text{ kg}$}
{$50\text{ kg}$}
{$54\text{ kg}$}
{$56\text{ kg}$}
\loigiai{
$80 - 28 = 52\text{ kg}$.\\
\textbf{Đáp án của câu hỏi là A}
}
\end{ex}

\begin{ex} %[2H7V35]
Một bể cá chứa $60\text{ l}$ nước, sau khi xả bớt $25\text{ l}$ thì trong bể còn lại số lít nước là
\choice
{$34\text{ l}$}
{$36\text{ l}$}
{$37\text{ l}$}
{\True $35\text{ l}$}
\loigiai{
$60 - 25 = 35\text{ l}$.\\
\textbf{Đáp án của câu hỏi là D}
}
\end{ex}

\begin{ex} %[2H7C35]
Thùng thứ nhất có $28\text{ l}$ dầu. Nếu đổ thêm $6\text{ l}$ vào thùng thứ nhất thì thùng thứ nhất vẫn ít hơn thùng thứ hai là $4\text{ l}$ dầu. Hỏi thùng thứ hai có bao nhiêu lít dầu?
\choice
{$36\text{ l}$}
{$40\text{ l}$}
{\True $38\text{ l}$}
{$34\text{ l}$}
\loigiai{
Thùng 1 sau khi thêm: $28 + 6 = 34\text{ l}$.\\
Thùng 2 có: $34 + 4 = 38\text{ l}$.\\
\textbf{Đáp án của câu hỏi là C}
}
\end{ex}

\textbf{Phần 2. Trắc nghiệm đúng sai}
\setcounterpageref{ex}{0}

\begin{ex} %[2H7H35]
Xét tính đúng sai về các đơn vị đo đại lượng:
\choiceTF
{\True $25\text{ kg} + 35\text{ kg} = 60\text{ kg}$.}
{\True $40\text{ l} - 15\text{ l} = 25\text{ l}$.}
{1 ngày có 12 giờ.}
{\True 1 ngày có 24 giờ.}
\loigiai{
\begin{itemize}
    \item $25 + 35 = 60\text{ kg}$: Đúng.
    \item $40 - 15 = 25\text{ l}$: Đúng.
    \item 1 ngày = 24 giờ (không phải 12 giờ). Ý 3 Sai, Ý 4 Đúng.
\end{itemize}
\begin{itemize}
    \item a) \textbf{Đúng}.
    \item b) \textbf{Đúng}.
    \item c) \textbf{Sai}.
    \item d) \textbf{Đúng}.
\end{itemize}
}
\end{ex}

\begin{ex} %[2H7V35]
Một can dầu chứa $18\text{ l}$, mẹ mua thêm $7\text{ l}$ dầu nữa. Xét tính đúng sai:
\choiceTF
{\True Tổng lượng dầu mẹ có là $25\text{ l}$ ($18 + 7 = 25$).}
{\True Nếu dùng hết $10\text{ l}$ dầu thì còn lại $15\text{ l}$ dầu ($25 - 10 = 15$).}
{Lượng dầu còn lại ít hơn $10\text{ l}$.}
{\True Lượng dầu còn lại là $15\text{ l}$ (nhiều hơn $10\text{ l}$).}
\loigiai{
\begin{itemize}
    \item Tổng: $18 + 7 = 25\text{ l}$: Đúng.
    \item Còn lại: $25 - 10 = 15\text{ l}$: Đúng.
    \item $15 > 10$ (không ít hơn). Ý 3 Sai, Ý 4 Đúng.
\end{itemize}
\begin{itemize}
    \item a) \textbf{Đúng}.
    \item b) \textbf{Đúng}.
    \item c) \textbf{Sai}.
    \item d) \textbf{Đúng}.
\end{itemize}
}
\end{ex}

\textbf{Phần 3. Trả lời ngắn}
\setcounterpageref{ex}{0}

\begin{ex} %[2H7H35]
Tính: $37\text{ kg} + 45\text{ kg}$ theo ki-lô-gam.
\shortans{82}
\loigiai{
\textbf{Đáp số:} $82$\\
$37 + 45 = 82\text{ kg}$.
}
\end{ex}

\begin{ex} %[2H7H35]
Tính: $65\text{ l} - 28\text{ l}$ theo lít.
\shortans{37}
\loigiai{
\textbf{Đáp số:} $37$\\
$65 - 28 = 37\text{ l}$.
}
\end{ex}

\begin{ex} %[2H7V35]
Bao gạo to nặng $50\text{ kg}$, bao gạo nhỏ nặng $25\text{ kg}$. Bao gạo to nặng hơn bao nhỏ bao nhiêu ki-lô-gam?
\shortans{25}
\loigiai{
\textbf{Đáp số:} $25$\\
$50 - 25 = 25\text{ kg}$.
}
\end{ex}

\begin{ex} %[2H7V35]
Một thùng có $45\text{ l}$ nước mắm, sau khi bán còn lại $19\text{ l}$. Hỏi đã bán bao nhiêu lít nước mắm?
\shortans{26}
\loigiai{
\textbf{Đáp số:} $26$\\
$45 - 19 = 26\text{ l}$.
}
\end{ex}

\begin{ex} %[2H7V35]
Tháng 12 có bao nhiêu ngày?
\shortans{31}
\loigiai{
\textbf{Đáp số:} $31$\\
31 ngày.
}
\end{ex}

\begin{ex} %[2H7C35]
Một cửa hàng buổi sáng bán được $35\text{ kg}$ gạo, buổi chiều bán được nhiều hơn buổi sáng $18\text{ kg}$ gạo. Hỏi cả ngày cửa hàng bán được bao nhiêu ki-lô-gam gạo?
\shortans{88}
\loigiai{
\textbf{Đáp số:} $88$\\
Buổi chiều: $35 + 18 = 53\text{ kg}$.\\
Cả ngày: $35 + 53 = 88\text{ kg}$.
}
\end{ex}
